\chapter{Optimization}
\section{Analysis on Covex Sets}
\subsection{Convex Set}
Recall the definition of linear combination:
$$\forall x,y \in V, a,b \in \mathbb{F}: ax+by.$$
If we further restrict $a+b=1$, we say it's a convex combination. 
\begin{definition}[Convex Set]
    Let $W$ be a subset of $V$, we say it's a convex subset if $W$ is closed under convex combination. We denote convex hull, the smallest convex set containing $W$, as $co(W)$, i.e.,
    $$co(W):=\bigcap\{C: C \text{ is convex } \land W\subset C\}.$$ 
\end{definition}
\subsection{Convex Functions}
\begin{definition}[Convex and Concave Functions]
    Suppose $C$ is a convex set, we say $f: C\to V$ to $V$ is (strictly) convex if and only if 
$\forall x,y\in C$ and $a\in [0,1]$, $$f(ax+(1-a)y) (<)\leq af(x)+(1-a)f(y).$$
\,If $-f$ is (strictly) convex, we say $f$ is (strictly) concave.
\end{definition}

We now turn studying whether a function is convex to studying whether a set is convex.
\begin{definition}
    We define the following:
    \begin{description}
        \item[Epigraph] $Epi(a):=\{(x,a) \in X\times \R: a\geq f(x)\}$
        \item[Subgraph] $Sub(a):=\{(x,a) \in X\times \R: a\leq f(x)\}$
    \end{description}
\end{definition}

\begin{proposition}
    Suppose $f:C\to \R$, then 
    \begin{itemize}
        \item $f$ is convex $\iff$ epigraph is convex;
        \item $f$ is concave $\iff$ subgraph is convex.
    \end{itemize}
\end{proposition}
\subsection{Quasi-Convexity}
We motivate the definition of quasi-convexity by the next example.
\begin{example}
    Consider a special Cobb-Douglas preference whose utility representation is $u(x,y) = x^{\frac{1}{2}}y^{\frac{1}{2}}$. If we add a monotonic transformation
to the utility representation, the new function is still an utility representation for the original preference. However, the concavity may NOT be preserved. For example, if we do the 
transformation $f(x) = x^6$, then $f\circ u (x,y)=x^3y^3$, which is convex, as you may verify. This leads to no difference in the UMP, but we want a concept that preserves concavity and also
leads to good properties when it comes to optimization.
\end{example}

\begin{definition}[Quasi-Convex and Quasi-Concave Functions]
    Suppose $C$ is a convex set, we say $f: C\to V$ to $V$ is (strictly) quasi-convex if and only if 
$\forall x,y\in C$ and $a\in [0,1]$, $$f(ax+(1-a)y) (<)\leq \max\{f(x),f(y)\}.$$
\,If $-f$ is (strictly) quasi-convex, we say $f$ is (strictly) quasi-concave.
\end{definition}

\begin{proposition}
    Suppose $g:\R\to \R$ is an increasing and $f: C\to \R$ is quasi-convex(quasi-concave), then $g\circ f$ is quasi-convex(quasi-concave).
\end{proposition}

We similarly define the concept of upper and lower contour set.
\begin{definition}[Upper and Lower Contour Set]
        We define the following:
    \begin{description}
        \item[Upper Contour Set] $U_f(a):=\{x\in X: f(x)\geq a\}$
        \item[Lower Contour Set] $L_f(a):=\{x \in X: f(x)\leq a\}$
    \end{description}
\end{definition}

\begin{proposition}
    Suppose $f:C\to \R$, then 
    \begin{itemize}
        \item $f$ is convex $\iff$ every lower contour set is convex;
        \item $f$ is concave $\iff$ every upper contour set is convex.
    \end{itemize}
\end{proposition}

\subsection{Application: Convexity of Preference and its Induced Utility Representation}


\subsection{Hyperplane Separation Theorem}
\begin{theorem}
    Suppose $C\subset \R^n$ is (closed and) convex, then there exists $p\neq 0 \in \R^n$ such that $$\sup_{y\in C} p\cdot y (<) \leq p\cdot x. $$
\end{theorem}

\begin{theorem}
    Suppose $A,B\subset \R^n$ are convex sets such that $A\cap B = \emptyset$, then there exists $p\neq 0 \in \R^n$ such that $$\sup_{x\in A} p\cdot x \leq \inf_{y\in B} p\cdot y. $$

    Moreover, if $A,B$ are closed and at least one is compact, then there exists $p\neq 0 \in \R^n$ such that $$\sup_{x\in A} p\cdot x < \inf_{y\in B} p\cdot y. $$
\end{theorem}

\subsection{Application: Strictly Dominated Stragtegy and Non Best Response}